\subsection{The exact zero-Euler classification at equal canonical order}

\begin{statement}{C.2.UN.65}{Proved}{Trivial local canonical class, triple points and the complete zero case}
Let $f:X\to\Delta$ be a proper degeneration of connected complex
two-tori, with smooth three-dimensional total space and full
reduced simple-normal-crossing central fibre
$S=\sum_iV_i$. Suppose that the orders of the actual canonical
evaluation along all $V_i$ are equal. Equivalently, the canonical
line of a neighbourhood of $S$ is trivial. Let $t$ be the number
of triple points and $r$ the number of rational double curves.
Then
\[
                         e(S)=t,\qquad 2r=3t.
\tag{C.CZ.1}
\]
In particular $e(S)$ is a nonnegative even integer.

The complete zero-Euler list in this class is as follows.
The fibre is either a smooth two-torus, or is a cycle of smooth
unblown components, each a marked diagonal primary or secondary
Hopf surface, or a marked surface
\[
                   \PP(\OO_E\oplus L),
          \qquad E\text{ elliptic},\quad L\in\operatorname{Pic}^0(E).
\tag{C.CZ.2}
\]
The boundary curves are exactly the two Hopf axes or the two
ruled sections. Their actual normal lines are dual across
every edge, and the actual integral complement-lattice product
is $M=I$.

Conversely every specified marked cycle in this list with
the stated dual-normal data and $M=I$ has a compact realization
over $\PP^1$, retaining all components and elliptic gluing
maps, with smooth total space and trivial canonical line
near the specified fibre. Its other singular fibre is the
conjugate marked cycle and all reduced fibre supports have
Euler number zero. The construction chooses normal-product
constants. A smooth two-torus is realized by its product
with $\PP^1$.
\end{statement}

\begin{proof}
\par\medskip\noindent\textbf{The volume form and the actual double-curve line maps.}
Shrink the disc to remove every other critical value.
The rank-one canonical direct image has an evaluation
generator $\Omega$ whose divisor is supported on $S$.
If all its orders are $k$, the quotient
$\Omega/f^k$ is a nowhere-zero holomorphic three-form.
This is the actual trivialization of $K_X$ used below.
Conversely a nowhere-zero three-form generates this
direct image: every other section is its product with
a holomorphic function constant on the compact connected
fibres, hence pulled back from the disc. Its evaluation
orders are all zero. This proves the equivalence in the
statement by the displayed multiplication and division maps.

Put $D_i=\sum_{j\ne i}(V_i\cap V_j)$, with each connected
double curve included once. Adjunction and the principal
divisor $S=(f)$ give the actual isomorphism
\[
                           K_{V_i}\simeq\OO_{V_i}(-D_i).
\tag{C.CZ.3}
\]
For a connected double curve $C\subset V_i\cap V_j$,
let $T_C$ be its reduced divisor of triple points.
Adjunction on $V_i$ gives
\[
             K_C=(K_{V_i}+C)|_C\simeq\OO_C(-T_C),
             \qquad 2g(C)-2=-\#T_C.
\tag{C.CZ.4}
\]
Every double curve is therefore elliptic and has no
triple points, or is rational and has exactly two.
Each triple point lies on exactly three double curves.
Counting these actual incidences proves $2r=3t$.

Restriction of the trivial line $\OO_X(S)$ to $C$
also gives the exact triple-point comparison
\[
 N_{C/V_i}\otimes N_{C/V_j}\otimes\OO_C(T_C)
                                  \simeq\OO_C.
\tag{C.CZ.5}
\]
Indeed $\OO_X(V_i)|_C=N_{C/V_j}$,
$\OO_X(V_j)|_C=N_{C/V_i}$, and every remaining
component contributes its actual transverse triple
point. In degrees this is
$C_{V_i}^2+C_{V_j}^2=-\#T_C$.
This is also the direct line-bundle proof of
\cite[Lemma 2.1]{Kulikov1977} in the present notation.

\par\medskip\noindent\textbf{The exact Euler identity.}
By (C.CZ.3), $K_{V_i}^2=D_i^2$.
The self-intersections of the double curves sum,
by (C.CZ.5), to $-3t$. The cross terms in the
boundary squares contribute $2$ on each of the
three surfaces through a triple point. Consequently
\[
                         \sum_iK_{V_i}^2=3t.
\tag{C.CZ.6}
\]
The alternating restriction sequence on the actual
SNC strata is
\[
 0\longrightarrow\OO_S\longrightarrow\bigoplus_i\OO_{V_i}
  \longrightarrow\bigoplus_C\OO_C
  \longrightarrow\bigoplus_p\C_p\longrightarrow0.
\tag{C.CZ.7}
\]
Proper flatness gives $\chi(\OO_S)=\chi(\OO_{F_t})=0$.
Only the rational double curves contribute to
$\sum_C\chi(\OO_C)$, so (C.CZ.7) gives
\[
                         \sum_i\chi(\OO_{V_i})=r-t.
\tag{C.CZ.8}
\]
The sequence and equality also agree with
\cite[Lemma 2.2]{Kulikov1977}.

Hirzebruch--Riemann--Roch and complex Gauss--Bonnet
on each compact complex surface give
$e(V_i)=12\chi(\OO_{V_i})-K_{V_i}^2$.
Topological inclusion-exclusion on the same strata
then yields
\[
 \begin{aligned}
 e(S)&=\sum_i e(V_i)-\sum_C e(C)+t\\
     &=12(r-t)-3t-2r+t\\
     &=10r-14t=t,
 \end{aligned}
\tag{C.CZ.9}
\]
where the last equality uses $2r=3t$.
These formulas require no K\"ahler assumption.
The count $2r=3t$ also makes $t$ even.

\par\medskip\noindent\textbf{The zero-case surfaces are minimal.}
Assume $e(S)=0$. Equations (C.CZ.1) and (C.CZ.4)
give no triple points and only elliptic double
curves. If there is just one component, the
central fibre is smooth and reduced. The
actual smooth-family comparison in UN5 makes
it a complex two-torus. We therefore treat
the connected reducible case, where every
$D_i$ is nonempty.

Equation (C.CZ.3) gives
$H^0(V_i,mK_{V_i})=H^0(V_i,\OO(-mD_i))=0$
for every $m>0$: its image in $H^0(V_i,\OO)=\C$
would be a constant vanishing on $D_i$.
Thus every component has Kodaira dimension
$-\infty$. The compact-surface classification
in C.3.1 makes its minimal model ruled or
of class VII.

For a ruled minimal model over a curve $B$
of genus at least two, every elliptic curve
on $V_i$ maps constantly to $B$, by
Riemann--Hurwitz. A general ruling fibre
$F\simeq\PP^1$ avoids their vertical images
and every blowup centre. It would satisfy
$D_i\cdot F=0$, whereas (C.CZ.3) and
adjunction give $D_i\cdot F=-K_{V_i}\cdot F=2$.
This contradiction leaves ruled genus zero
or one, and class VII.

Their coherent Euler characteristics are,
respectively, one, zero and zero. Equation
(C.CZ.8) is now $\sum_i\chi(\OO_{V_i})=0$,
so no rational component occurs. A ruled
component over an elliptic curve has Euler
number equal to its number of point blowups.
A class VII component has $b_1=b_3=1$ and
$e(V_i)=b_2(V_i)\ge0$. Since all double
curves are elliptic, (C.CZ.9) is also
$0=e(S)=\sum_i e(V_i)$. Every term is
therefore zero. The ruled components have
no blowups. Every class VII component is
minimal with $b_2=0$, and its actual
minimal-model restriction in UN19 makes
it a Hopf surface. This proves the claimed
absence of blowups and identifies both
remaining kinds of surface.

\par\medskip\noindent\textbf{The ruled boundary consists of two actual sections.}
Let $\pi:V\to E$ be one of the minimal ruled
surfaces over an elliptic curve. Each component
of $D$ is elliptic and maps finitely to $E$:
it cannot be contained in a rational ruling
fibre. The map is unramified by
Riemann--Hurwitz. Since $D\cdot F=2$, the
only possibilities are two degree-one
sections or one connected elliptic
bisection $C$.

The latter is excluded by an exact residue
map. Suppose $K_V+C\simeq\OO_V$.
The sequence
\[
 0\longrightarrow K_V\longrightarrow K_V(C)
      \xrightarrow{\operatorname{res}}K_C\longrightarrow0
\tag{C.CZ.10}
\]
has a nonzero residue of the nowhere-zero
section of $K_V(C)$: in a local equation
$x=0$ it is the residue of a unit times
$dx\wedge dz/x$. Thus
$H^0(K_V(C))\to H^0(K_C)$ is onto and
its connecting map to $H^1(K_V)$ is zero.
Under Serre duality that connecting map
is dual to the actual restriction
$H^1(V,\OO_V)\to H^1(C,\OO_C)$.
But $\pi^*:H^1(E,\OO_E)\to H^1(V,\OO_V)$
is an isomorphism, and its composite with
restriction is pullback by the degree-two
unramified map $C\to E$. The trace of
that pullback is multiplication by two,
so it is nonzero. This contradicts the
zero connecting map and rules out the
bisection.

The two disjoint sections give two
quotient lines of the rank-two ruling
bundle. Their combined quotient map is
an isomorphism on every fibre, because
the sections are distinct. It is therefore
an actual bundle isomorphism with the
sum of those two lines. Tensoring by
one inverse line gives
$V\simeq\PP(\OO_E\oplus L)$ with the
two specified sections. At this point
$L$ may have any degree; that degree
is determined below by the actual
nearby fundamental group.

\par\medskip\noindent\textbf{The Hopf boundary is the individually marked diagonal pair.}
Pull $K_H+D\simeq\OO_H$ back to an actual
central primary Hopf cover, as in UN43.
For a nonlinear contraction the only
elliptic curve is its axis $C$, with
$K=-(a+1)C$, $a\ge1$. The effective
axis line has infinite order. Thus a
reduced nonempty pullback boundary
cannot satisfy $K+D=0$ there.

For a diagonal contraction, the lifted
anticanonical section is
$g(z,w)\partial_z\wedge\partial_w$,
where Hartogs extends $g$ across zero
and
\[
                        g(\alpha z,\beta w)
                                      =\alpha\beta g(z,w).
\tag{C.CZ.11}
\]
In the nonresonant case its only
possible monomial is $zw$.
In the resonant case a further central
power has multipliers $(\gamma^p,\gamma^q)$,
with coprime positive $p,q$.
Its monomials have weight $p+q$.
For $p,q>1$ this again gives only
$zw$. For $p>1,q=1$ the equation is
$g=w(az+bw^p)$, and reducedness
requires $a\ne0$. The triangular
coordinate change $z\mapsto z+(b/a)w^p$
commutes with the contraction and
gives the two axes. The symmetric
case is identical. When $p=q=1$,
the reduced quadratic is a product
of two distinct linear forms, and
an invertible linear change commuting
with the scalar contraction gives
the axes. Thus the pulled boundary
is exactly two axes in every case.

Their unique logarithmic two-form,
up to a constant, is
$dz/z\wedge dw/w$. A deck lift
preserving both axes individually
is diagonal by the contraction-unit
argument of UN59 and fixes this
form. A lift interchanging the axes
is a diagonal map followed by the
interchange and negates the form.
The nonzero section of $K_H(D)$
pulls back to a deck-invariant
section, so no interchange occurs.
Consequently the two images are
distinct elliptic curves, the full
deck group is diagonal, and the
original marked pair is precisely
one of the diagonal primary or
secondary Hopf pairs of UN57.

\par\medskip\noindent\textbf{The integral obstruction to nonzero ruled degree.}
Every component now has exactly two
boundary curves. The connected
incidence graph is a cycle. At a
ruled vertex its complement is the
actual nonzero line $L^\times$ over
$E$. For $d=\deg L$ its group is
\[
 \pi_1(L^\times)=
 \langle a,b,c\mid [a,c]=[b,c]=1,
                               \ [a,b]=c^d\rangle .
\tag{C.CZ.12}
\]
This is the actual circle-bundle
extension: trivialize over an
oriented elliptic parallelogram;
the boundary commutator is the
degree-$d$ clutching loop of its
fibre. The fibre generator $c$
has infinite order, since the
homotopy sequence injects
$\pi_1(\C^*)$ when $\pi_2(E)=0$.

For a given smoothing take the
component cores and the actual
$xy=t$ collars. A Hermitian radial
retraction identifies a ruled
core with an annulus bundle in
$L^\times$; either boundary
inclusion induces the isomorphism
onto (C.CZ.12). For a diagonal
Hopf core, the same statement is
its actual rank-three logarithmic
lattice comparison. The dual-normal
maps identify the intervening
circle bundles. The canonical
divisor sections and a small smooth
unit root give these product-preserving
tubular maps for the arbitrary
given smoothing, exactly as in
UN46. Flatness of $L$ is not used
for this topological comparison.

The resulting cycle of group
isomorphisms is a graph of groups
whose vertex groups inject into
the total group. Explicitly cut
one edge, compose the remaining
isomorphisms, and obtain a single
vertex group with one stable
letter implementing the closing
automorphism. Its normal form
contains that vertex group as
a subgroup. Since the actual
nearby group is $\Z^4$, every
injected vertex group is abelian.
Equation (C.CZ.12) and the
infinite order of $c$ force
$d=0$ at every ruled vertex.
This proves (C.CZ.2) by the
actual collar and clutching maps.

The remaining component groups
are all their actual rank-three
abelian lattices. The same cycle
calculation gives
$\pi_1(F_t)=\Lambda\rtimes_M\Z$.
Its injected lattice and the
abelian group $\Z^4$ force
$M=I$. We have obtained exactly
the marked data of UN57, with
no omitted unblown component
or boundary alternative.

\par\medskip\noindent\textbf{Realization with a nowhere-zero local volume.}
For every such marked cycle,
UN57 gives the actual proper
plumbing, the compact two-ended
period family, and its exact
punctured-family gluing. It
retains all marked components
and elliptic maps. We check
the required local canonical
condition on this construction.

On its common period charts
put
\[
                       \Omega=dt\wedge d\eta_1\wedge d\eta_2.
\tag{C.CZ.13}
\]
Here $\eta$ is the common
vector coordinate obtained
from each component chart
by its actual moving frame
$P_i(t)$. Thus every overlap
has identity linear part in
this common frame.
Every period and affine edge
translation depends holomorphically
on $t$; its differential changes
the vector differentials only by
multiples of $dt$. Thus this
three-form descends through the
actual lattice and collar maps.
It is nonzero on component cores.
At a double curve the relative
two-form is a holomorphic unit
times $dx/x\wedge dz$, where
$z$ is an elliptic coordinate
and $xy=t$. Hence
\[
             dt\wedge(dx/x\wedge dz)
                               =dy\wedge dx\wedge dz.
\tag{C.CZ.14}
\]
The right side is holomorphic
and nowhere zero. The moving
Hopf and ruled frames have
nonzero determinant at the
central value, so the displayed
unit never vanishes after
shrinking the disc. The form
therefore extends across all
of $S$ as a nowhere-zero
three-form. Its evaluation
orders are equal to zero.
The compact realization in
the statement now follows
from UN57, with the same
verification at its conjugate
fibre. The smooth torus case
is its literal product family.
\end{proof}

\par\medskip\noindent\textbf{Exact relation to the primary modification literature.}
Kulikov's definition of a good degeneration in
\cite[p.~958]{Kulikov1977} is the smooth-total,
reduced SNC condition used here. Its Lemmas 2.1--2.2
give the same principal-divisor and coherent-stratum
maps as (C.CZ.5)--(C.CZ.8). The general canonical
modification statement is
\cite[Theorem 1$'$]{Kulikov1981}: a good degeneration
of surfaces with trivial canonical class admits a
good modification with trivial total canonical
class precisely under its stated minimum-order
condition on class VII components. The present
equal-order hypothesis gives the actual identity
comparison to such a target through $\Omega/f^k$.
The zero-case classification above is proved on
that given neighbourhood and supplies the explicit
marked-cycle maps and compact occurrence. The
1981 appendix also records the correction to the
1977 Enriques modification assertion; the cited
line and Euler comparisons are the displayed
1977 lemmas.

Kamenova's author-hosted primary text
\cite[Theorem 3.1]{KamenovaTori} assumes a
K\"ahler total space and trivial $K_X$.
Its necessary list contains a smooth torus,
elliptic-ruled paths and cycles, and a rational
case with a triangulated torus as dual complex.
The residue and finite-trace maps in
(C.CZ.10) exclude a terminal elliptic bisection,
so the path alternative contributes no
additional zero-Euler support in the stated
smooth-component setting. This sharpens that
necessary list. In the K\"ahler case Hopf
components are also excluded by restriction
of the K\"ahler form: their vanishing real
$H^2$ would make its positive square integrate
to zero by Stokes. The present proof keeps
the actual component and collar maps without
a K\"ahler hypothesis and supplies compact
occurrence of every marked zero-Euler cycle
specified in the statement.
